\documentclass[11pt]{article}
\usepackage[margin=2.5cm]{geometry}
\setlength{\parindent}{0pt}\setlength{\parskip}{0.8em}
\begin{document}
Hung Yen, Vietnam\\
\today

The Editors\\
\textit{Decision Analytics Journal}

\textbf{Re: Submission of ``Adaptive Influence Minimization via Deferred Frontier Blocking'' (manuscript DAJOUR-D-26-00787)}

Dear Editors,

We submit the above manuscript for consideration as a research article in \textit{Decision Analytics Journal}.

Blocking a rumor or epidemic that has already started, by removing a few nodes under a fixed budget, is usually treated as a one-shot decision made before the outbreak is observed. In practice a moderator or health authority sees where the cascade actually goes and can hold part of the budget back. The paper asks whether doing so can hurt, and how much it can help.

We introduce DEFER, a wrapper that turns any one-shot influence minimization planner into an adaptive one by re-planning every round and blocking a planned node only once the cascade has reached its neighborhood. Our main results are:
\begin{itemize}
\item pure deferral is never worse in expectation than the planner it wraps, subject to a held-out test for plan replacement whose residual error we do not bound;
\item a local rule for early commitment is optimal for a single target but can fail when deferred targets compete for the budget, which we show by counterexample;
\item the adaptivity gap of influence minimization is unbounded, with an explicit construction.
\end{itemize}
In experiments in which every policy faces the same random cascades, DEFER saves 36 to 97 percent of the spread on the family used in the proof and beats greedy, re-planning and one-shot integer-programming policies in all 18 settings. On 45 synthetic settings it is significantly better than greedy in 9 and never significantly worse. On four empirical and synthetic networks it is significantly better in 5 of 40 settings, with small typical differences. Integer programs over sampled scenarios did not beat greedy or DEFER and ran far slower. We report the settings in which adapting does not help as well as those in which it does.

We believe the paper fits the journal's readership because it gives decision makers in misinformation control and epidemic response a concrete rule for when to spend an intervention budget at once and when to wait, together with guarantees and the computational cost of acting on it.

The manuscript is original, is not under consideration elsewhere, and has been approved by all authors. The authors declare no competing interests. A title page with author details, a blinded manuscript and the highlights are enclosed.

Thank you for your consideration.

Sincerely,

Quang-Vinh Dang (on behalf of all authors)\\
British University Vietnam, Hung Yen, Vietnam\\
vinh.dq4@buv.edu.vn
\end{document}
